\documentclass[11pt]{article}
\usepackage[T1]{fontenc}
\usepackage{lmodern,amsmath,amssymb,amsthm,mathtools,mathrsfs,microtype}
\usepackage[margin=1.03in]{geometry}
\usepackage[hidelinks]{hyperref}
\hypersetup{pdftitle={The generalized Mukai conjecture},pdfauthor={OpenAI}}
\newtheorem{theorem}{Theorem}[section]
\newtheorem{proposition}[theorem]{Proposition}
\newtheorem{lemma}[theorem]{Lemma}
\newtheorem{corollary}[theorem]{Corollary}
\theoremstyle{definition}
\newtheorem{definition}[theorem]{Definition}
\theoremstyle{remark}
\newtheorem{remark}[theorem]{Remark}
\DeclareMathOperator{\Tr}{Tr}
\DeclareMathOperator{\NS}{NS}
\DeclareMathOperator{\ev}{ev}
\DeclareMathOperator{\Pic}{Pic}
\DeclareMathOperator{\Mor}{Mor}
\newcommand{\pt}{\mathrm{pt}}
\newcommand{\dd}{\mathrm d}
\newcommand{\C}{\mathbb C}
\newcommand{\Pj}{\mathbb P}
\pdfinfoomitdate=1
\pdfsuppressptexinfo=15
\pdftrailerid{}
\title{The generalized Mukai conjecture}
\author{OpenAI}
\date{September 24, 2026}
\begin{document}
\maketitle
\begin{abstract}
We prove the generalized Mukai conjecture: every positive-dimensional
smooth complex Fano manifold of dimension $n$, Picard number $\rho$, and
pseudoindex $\iota$ satisfies $\rho(\iota-1)\le n$. Equality holds
precisely for the product of $\rho$ copies of $\mathbb P^{\iota-1}$.
\end{abstract}
\tableofcontents
\clearpage
\section{Introduction}\label{sec:introduction}

A Fano manifold is a smooth connected complex projective variety $X$
whose anticanonical divisor $-K_X$ is ample. Its Picard number
$\rho_X$ is the rank of its N\'eron--Severi group, or equivalently the
dimension of the space of numerical divisor classes. If $X$ has positive
dimension, its \emph{pseudoindex} is
\[
 \iota_X=\min\{-K_X\cdot C:C\subset X
                  \text{ is an irreducible rational curve}\}.
\]
The existence of rational curves on a Fano manifold makes this minimum
well defined. The pseudoindex records the smallest anticanonical curve
degree; it need not equal the largest integer dividing $-K_X$ in
$\operatorname{Pic}(X)$, which is the Fano index.

The generalized Mukai conjecture asks for a sharp bound on the Picard
number in terms of dimension and pseudoindex, together with a
classification of the extremal manifolds. We prove the conjectured bound
and classification in arbitrary dimension.

\begin{theorem}\label{thm:main}
Let $X$ be a smooth connected complex projective Fano manifold of
positive dimension $n$, Picard number $\rho_X$ and pseudoindex $\iota_X$.
Then
\[
 \rho_X(\iota_X-1)\le n.
\]
Equality holds if and only if
\[
 X\cong(\mathbb P^{\iota_X-1})^{\rho_X}.
\]
\end{theorem}

The index version originates in Mukai's work~\cite{Mukai1988}.
Wi\'{s}niewski introduced the pseudoindex and proved that
$2\iota_X>n+2$ forces $\rho_X=1$~\cite[Theorem~A]{Wisniewski1990}.
Bonavero, Casagrande, Debarre and Druel formulate the pseudoindex
strengthening and establish it for Fano manifolds of dimension at most
four, and for smooth toric Fano manifolds when $n\le7$ or
$3\iota_X\ge n+3$~\cite{BCDD2003}. Andreatta, Chierici and Occhetta
extend the low-dimensional result to dimension five and prove
high-pseudoindex cases subject to contraction hypotheses
\cite{AndreattaChiericiOcchetta2004}.

For varieties with additional symmetry,
Casagrande settles the bound and equality classification for
$\mathbb Q$-factorial Gorenstein toric Fano varieties in every
dimension~\cite[Theorem~1(ii)]{Casagrande2006}, and Pasquier treats
$\mathbb Q$-factorial Gorenstein horospherical Fano
varieties~\cite[Theorem~1]{Pasquier2010}. More recently, Gagliardi,
Hofscheier and Pearson prove the generalized conjecture for
$\mathbb Q$-factorial spherical Fano varieties
\cite{GagliardiHofscheierPearson2025}.

Large pseudoindex gives a different route. Novelli and Occhetta
prove both conclusions for all smooth Fano manifolds with
$\iota_X\ge(n+3)/3$~\cite{NovelliOcchetta2010}, and Novelli later
enlarges this range to $\iota_X>n/3$~\cite[Theorem~5.1]{Novelli2012}.
At the boundary $n=3\iota_X$, Novelli also proves the conjecture
when $X$ admits an extremal contraction of fiber type
\cite{Novelli2026}. Theorem~\ref{thm:main}
imposes no restriction on the dimension, Picard number or geometry
of the smooth Fano manifold.

The ordered-chain method of Bonavero, Casagrande, Debarre and Druel
derives the inequality once $\rho_X$ numerically independent
proper irreducible components of the space of rational curves, joined
by a nonempty chain, are available~\cite[Corollary~5.3]{BCDD2003}.
Our argument first detects
independent curve degrees in quantum correspondences. Proper covering
families enter only at equality, when all the selected degrees are minimal.

\subsection*{The argument}

Write $r=\rho_X$ and $\iota=\iota_X$. If $\iota=1$, the inequality
is strict because $n>0$, so assume $\iota\ge2$.
Choose divisor classes $D_1,\ldots,D_r$ forming a numerical basis.
Small quantum multiplication by these divisors gives commuting
matrices $A_1(q),\ldots,A_r(q)$ on the even cohomology of $X$,
with coefficients depending on $r$ nonzero formal variables $q_j$.
Each matrix consists of classical cup product and a finite sum of
two-point curve correspondences. A correspondence of anticanonical
degree $d$ lowers cohomological codimension by $d-1$.
Consequently, a nonzero product in $r$ independent curve degrees would
give
\[
 r(\iota-1)\le\sum_{j=1}^r(d_j-1)\le n.
\]
The main task is to produce that product without assuming that quantum
cohomology is semisimple.

We first prove a lower bound for a point descendant, an intersection
number on the space of one-pointed rational stable maps involving the
$(d-2)$nd power of the marked cotangent class. Here $d$ is the
anticanonical degree. For any degree represented by a free rational map
(a map whose pulled-back tangent bundle is globally generated), this
number is at least $(ad)^{-(d-1)}$, where $-aK_X$
is a fixed very ample divisor. To prove positivity, we contract all
components away from the marking and map the evaluation fiber to a
weighted projective stack. The image of its boundary has dimension
too small to contribute. The pushed-forward virtual class is therefore
a nonempty effective cycle, whose degree admits the required uniform
bound. Section~\ref{sec:descendants} gives the construction in families
and keeps track of stack multiplicities.

Two-point descendant solutions of quantum differential equations
provide a conceptual antecedent for the recurrence used next
\cite[\S6]{Givental1996}.
Here the divisor and recursion identities give a finite-shift
recurrence for descendant vectors on an arbitrary smooth Fano manifold.
After factorial normalization it yields an
exponential upper bound. If every joint eigenvalue tuple of the quantum
divisors satisfied a polynomial relation, the same recurrence would
force decay of order $\exp(-c d\log d)$ along a suitable free curve
degree. After the same factorial normalization, the descendant lower
bound permits only exponential decay and rules this out. Thus one joint
tuple has $r$ algebraically independent coordinates
(Section~\ref{sec:spectral}).

Section~\ref{sec:trace} converts this algebraic independence into the
required curve correspondences. An alternating trace of matrix-valued
differentials is invariant even under changes of basis depending on
$q$. Computing it in a common triangular basis detects the independent
eigenvalues; computing it in the fixed cohomology basis detects a
nonzero product in independent curve degrees. This argument also
applies to nonsemisimple commuting matrices.

At equality all these degrees are minimal. Their nonzero product
implies the existence of a chain of actual rational curves; minimality
makes the full irreducible families containing them proper.
The ordered-chain bound of Bonavero, Casagrande, Debarre and Druel
applies to these components, as required by its erratum
\cite{BCDDErratum}. A cyclic reordering makes every family covering,
and Occhetta's characterization of products of projective spaces gives
the asserted rigidity~\cite{Occhetta2006}.
Section~\ref{sec:equality} verifies each of these family hypotheses.

\section{Quantum divisors and point descendants}\label{sec:setup}

We set up a finite collection of commuting matrices and a sequence of
vectors on which they act. Their grading records anticanonical degrees;
the vectors will connect their joint eigenvalues to rational curves.
Throughout the proof, $X$ is a smooth projective complex Fano manifold,
$n=\dim X>0$, $r=\rho_X$, and $\iota=\iota_X\ge2$.
The case $\iota=1$ will require no argument beyond $n>0$.

Choose an integral basis $D_1,\ldots,D_r$ of the N\'eron--Severi group
modulo torsion. A numerical curve degree is represented by
\[
 \beta=(D_1\cdot\beta,\ldots,D_r\cdot\beta)\in\mathbb Z^r.
\]
Here and below we use the same symbol for a numerical degree and its
coordinate vector. Write
\[
 -K_X\equiv\sum_{j=1}^r c_jD_j,
 \qquad d_\beta=\sum_{j=1}^r c_j\beta_j.
\]
A positive degree means the degree of a nonconstant genus-zero stable
map. Every such degree has $d_\beta\ge\iota$. There are only finitely
many positive numerical degrees with bounded $d_\beta$: for a
sufficiently large integer $a$, both $a(-K_X)+D_j$ and $a(-K_X)-D_j$
are ample, so $|\beta_j|\le a d_\beta$ on effective degrees.
When several curve classes have the same numerical degree, we sum their
invariants. Proper finite-type spaces of stable maps of bounded
projective degree make these sums finite.

We use ordinary genus-zero Gromov--Witten invariants, defined by the
virtual fundamental classes of stable-map stacks. An insertion
$\tau_\ell(a)$ means $\psi^\ell\ev^*a$ at the indicated marking:
$\ev$ evaluates the stable map there, and $\psi$ is the first Chern
class of its cotangent line.
The divisor, string and topological recursion equations, together with
associativity, are the standard identities we use
\cite{BehrendVirtual,KontsevichManinQuantum,KontsevichManinDescendants,PandharipandeDescendants}.
The discussion below specifies all unstable conventions needed here.
Set
\[
 H=\bigoplus_{k=0}^n H_k,
 \qquad H_k=H^{2k}(X,\mathbb C),
 \qquad (a,b)=\int_Xa\cup b,
\]
and put $H_k=0$ outside this range. The pairing on $H$ is nondegenerate.
We work only with even cohomology; invariants involving a single odd
intermediate class and otherwise even insertions vanish by parity.

For a positive numerical degree $\beta$, define $S_\beta\in\operatorname{End}(H)$ by
\[
 (S_\beta a,b)=\langle a,b\rangle_\beta.
\]
The virtual dimension of the two-pointed space is $n+d_\beta-1$.
Consequently
\begin{equation}\label{eq:grading}
 S_\beta(H_k)\subset H_{k+1-d_\beta}.
\end{equation}
In particular $S_\beta=0$ for $d_\beta>n+1$. For nonzero formal variables
$q_1,\ldots,q_r$, write $q^\beta=\prod_jq_j^{\beta_j}$ and define
\begin{equation}\label{eq:quantum-operators}
 A_j(q)=D_j\cup(-)+\sum_{\beta>0}q^\beta\beta_jS_\beta
       =\sum_\gamma q^\gamma A_{j,\gamma}.
\end{equation}
Thus $A_{j,0}=D_j\cup(-)$, $A_{j,\beta}=\beta_jS_\beta$ for positive
degrees, and all other coefficients vanish. The grading bound and
finiteness of bounded degrees make these Laurent polynomial matrices.
The divisor equation identifies them with small quantum multiplication
by $D_j$, so associativity implies that they commute. If $R_t$ acts
by $t^k$ on $H_k$, then
\begin{equation}\label{eq:homogeneity}
 A_j(t^{c_1}q_1,\ldots,t^{c_r}q_r)
       =tR_t^{-1}A_j(q)R_t.
\end{equation}
Indeed the classical term raises the grading by one, whereas the
$\beta$ term changes it by $1-d_\beta$.

\subsection{The descendant recurrence}

Two-point descendant generating series also occur in Givental's
quantum differential framework~\cite[\S6]{Givental1996}.  The
recurrence below is derived directly from the virtual genus-zero
identities cited above, so it does not require a convexity assumption
on $X$.

Let $\pt\in H_n$ be the class of a point, normalized by
$\int_X\pt=1$. Define $v_0=\pt$ and, for positive degrees, define
$v_\beta\in H$ by
\[
 (v_\beta,b)=\sum_{\ell\ge0}
             \langle\tau_\ell(\pt),b\rangle_\beta.
\]
The dimension constraint makes each sum finite. Set all other
$v_\beta$ equal to zero. In particular there is no degree-zero
unstable two-point invariant in this definition.

\begin{lemma}\label{lem:recurrence}
For every integer degree vector $\beta$ with $d_\beta>0$ and every
$j\in\{1,\ldots,r\}$,
\begin{equation}\label{eq:recurrence}
 \beta_jv_\beta=\sum_\gamma A_{j,\gamma}v_{\beta-\gamma}.
\end{equation}
For a positive degree $\beta$,
\begin{equation}\label{eq:string-pairing}
 (v_\beta,1)=\langle\tau_{d_\beta-2}(\pt)\rangle_\beta.
\end{equation}
\end{lemma}

\begin{proof}
Pair the first identity with an arbitrary even class $b$. Inserting
$D_j$ into $\langle\tau_\ell(\pt),b\rangle_\beta$ multiplies the invariant
by $\beta_j$ \cite[Lemma~1.4, equation~(12)]{KontsevichManinDescendants}: the descendant correction contains
$D_j\cup\pt=0$. For $\ell>0$, genus-zero topological recursion at the
point marking \cite[equation~(4a)]{KontsevichManinDescendants} separates a two-point factor containing
$\tau_{\ell-1}(\pt)$ from a three-point factor containing $D_j$ and $b$.
The point side has positive degree, since a two-point degree-zero
component is unstable. The other side can have degree zero, in which
case it contributes the classical cup operator. Summing over $\ell>0$
therefore gives the terms with positive-degree $v_{\beta-\gamma}$.
The $\ell=0$ primary invariant gives the remaining term
$A_{j,\beta}v_0$. Splitting via the Poincar\'e pairing gives exactly
\eqref{eq:recurrence}. If a degree is not represented, the same splitting
identity holds with zero coefficients.

For the second identity, dimension leaves only $\ell=d_\beta-1$ in the
pairing with $1$. The string equation \cite[\S1.2, p.~311]{PandharipandeDescendants} then gives
\[
 \langle\tau_{d_\beta-1}(\pt),1\rangle_\beta
 =\langle\tau_{d_\beta-2}(\pt)\rangle_\beta.
\]
The exponent is nonnegative because $d_\beta\ge\iota\ge2$.
\end{proof}

The next section gives a positive lower bound for this pairing whenever
$\beta$ has a free rational representative. It is a statement about a
pushforward of the virtual evaluation fiber; it will not require the
entire virtual class to be effective.

\section{A lower bound for point descendants}
\label{sec:descendants}

The spectral argument will require point descendants in a sequence of
large curve degrees.  Their virtual definition does not give positivity
directly.  We prove the required lower bound by a contraction that forgets
all components away from the marked component.  Its boundary image is too
small to contribute to the integral.

Throughout this section, \(X\) is a smooth projective complex Fano manifold
of dimension \(n>0\) and pseudoindex at least two.  Fix an embedding
\[
  \jmath:X\hookrightarrow\mathbb P^N,\qquad
  \jmath^*\mathcal O_{\mathbb P^N}(1)\cong\mathcal O_X(-aK_X),
\]
where \(a\) is a positive integer.  A morphism
\(f:\mathbb P^1\to X\) is \emph{free} if \(f^*T_X\) is globally generated.
For a numerical curve degree \(\beta\), set \(d_\beta=-K_X\cdot\beta\).
The notation \(\overline{\mathcal M}_{0,1}(X,\beta)\) includes the union
over actual curve classes with numerical degree \(\beta\).  This is a
finite union: maps of fixed projective degree belong to a finite-type
proper moduli stack, and numerical degree is locally constant.

We use rational Chow groups of Deligne--Mumford stacks.  In particular,
degrees of zero-dimensional stack cycles include their stabilizer
denominators.  The marked cotangent line has first Chern class \(\psi\).

\begin{proposition}[Point-descendant lower bound]
\label{prop:descendant-bound}
Suppose that \(\beta\) is represented by a nonconstant free map
\(\mathbb P^1\to X\).  Put \(d=d_\beta\) and \(b=ad\).  Then
\[
  \left\langle\tau_{d-2}(\mathrm{pt})\right\rangle_\beta
  \ge \frac{1}{b^{d-1}}.
\]
\end{proposition}

We will express this descendant as the degree of a nonempty effective
cycle in a weighted projective stack whose weights are at most \(b\).
The main task is to obtain that cycle from the virtual evaluation
fiber: irreducible maps will contribute with stack degree one, while
the boundary image will have too small a dimension to contribute.
After proving the bound, we show that free degrees are Zariski dense
in the numerical curve space.

\subsection{A free evaluation fiber}

There is a very general point \(x\in X\) such that every nonconstant
morphism \(\mathbb P^1\to X\) whose image contains \(x\) is free.
Here is the argument, including the quantifier over special maps.
Stratify the nonfree locus in each morphism scheme into smooth
irreducible locally closed reduced strata.  There are countably many
such strata.  If the evaluation from \(\mathbb P^1\times T\), for one
stratum \(T\), were dominant, generic smoothness would give a point
\((p,f)\) where its differential is surjective.  Its image is contained
in the image of
\[
  H^0(\mathbb P^1,f^*T_X)\longrightarrow(f^*T_X)_p.
\]
Indeed, moving \(p\) gives directions also obtained from
reparametrizations of the source.  Splitting
\(f^*T_X=\bigoplus_j\mathcal O(a_j)\), surjectivity at one point implies
\(a_j\ge0\) for every \(j\), contrary to the choice of \(T\).
Thus each such evaluation image has proper closure, and their
complement gives the assertion.

Evaluation is smooth at a free map: the obstruction groups
\(H^1(\mathbb P^1,f^*T_X)\) and
\(H^1(\mathbb P^1,f^*T_X(-p))\) vanish.  Consequently the degree
\(\beta\) in Proposition~\ref{prop:descendant-bound} has maps
through an open subset of \(X\).  Choose \(x\) in this open subset and
outside all the exceptional closed subsets just considered.

Let
\[
  F=\operatorname{ev}^{-1}(x)
    \subset\overline{\mathcal M}_{0,1}(X,\beta),\qquad h=d-2.
\]
The fiber carries the refined virtual class
\[
  [F]^{\mathrm{vir}}
  =x^![\overline{\mathcal M}_{0,1}(X,\beta)]^{\mathrm{vir}}
  \in A_h(F)_{\mathbb Q}.
\]
Let \(U\subset F\) be the open substack with irreducible source.
It is nonempty, smooth of dimension \(h\), and
\[
  [F]^{\mathrm{vir}}|_U=[U].
\]
Only this open part is asserted to be unobstructed.

To use this equality of classes on \(U\), we will construct a proper
morphism that preserves the marked cotangent class, has generic stack
degree one on \(U\), and sends \(F\setminus U\) to a locus of
dimension less than \(h\).  These properties will make the
pushforward of \([F]^{\mathrm{vir}}\) effective.

\subsection{Contracting away from the marking}

Compose maps with \(\jmath\).  Their projective degree is \(b\), and
this composition changes no stable domain, since it contracts no
nonconstant component.  At a geometric point, retain the component
carrying the marking and replace each attached tree by a common zero
of the map sections at its attachment point.  More precisely, multiply
the original sections restricted to the marked component by a common
factor whose vanishing order at each attachment is the projective
degree of the removed tree.  The resulting tuple has total degree
\(b\) and no common zero at the marking.  This description also
allows the marked component to be constant.

We now construct these data in families on
\(\overline{\mathcal M}_{0,1}(\mathbb P^N,b)\) and track the marked
cotangent line.  The contraction also follows from the one-marked
intermediate-space construction of
Mustaţă--Mustaţă~\cite[Definition~1.2 and Proposition~1.3]
{MustataMustata2007}, with the fine-stack construction of
\cite[Proposition~1.7]{MustataMustataChow}.

Write \(\pi:\mathcal C\to\overline{\mathcal M}_{0,1}(\mathbb P^N,b)\)
for the universal curve, \(s\) for its marking, and
\(\mathcal L\) for the pullback of \(\mathcal O_{\mathbb P^N}(1)\).
The moduli stack is smooth and its boundary has the usual independent
node-smoothing parameters: on a genus-zero tree,
\(H^1(C,f^*T_{\mathbb P^N})=0\).  Thus the universal curve is regular,
with local nodal models \(uv=t\).

For each boundary branch separating an unmarked tail of total degree
\(e\), take in a local node-smoothing chart the Cartier divisor
consisting of its unmarked side.  The marked side and degree \(e\)
are intrinsic, so summing these local divisors over all such branches
is invariant under permutations of equal-degree branches and glues
in families.  Let \(\Delta_e\) be this sum, counting repeated branches
once each.
Only \(1\le e\le b-1\) occurs.  A tail of degree \(b\) would leave a
degree-zero marked side with only one external attachment; a finite
tree of contracted components with one marking and one external
attachment cannot satisfy stability.

Set
\[
  \mathcal L^+
    =\mathcal L\Bigl(\sum_{e=1}^{b-1}e\Delta_e\Bigr).
\]
On a nodal fiber, twisting by an unmarked-side divisor transfers one
unit of degree across that node toward the marked side.  To compute
the resulting multidegree, root the dual tree at the marked component.
For a nonroot component \(v\), let \(b_v\) be its original degree
and \(E_v\) the total degree of the subtree rooted at \(v\).
If \(w\) runs over its children, then
\[
  \deg_v\mathcal L^+=b_v-E_v+\sum_w E_w=0.
\]
On the root, of original degree \(b_0\), the degree is
\(b_0+\sum_wE_w=b\).  Thus \(\mathcal L^+\) has degree \(b\)
on the component containing \(s\), and degree zero elsewhere.

The line bundle \(\mathcal O_{\mathcal C}(s)\) has degree one on that
marked component and degree zero elsewhere.  On each fiber it is
globally generated with \(h^0=2\) and \(h^1=0\).  Cohomology and base
change consequently give a rank-two bundle
\(\pi_*\mathcal O_{\mathcal C}(s)\), and its evaluation map contracts
the universal curve onto a \(\mathbb P^1\)-bundle.  This contraction
is an isomorphism on the marked component and near the section.

The line bundle \(\mathcal L^+(-bs)\) has multidegree zero on every
fiber.  Its fiberwise \(h^0=1\), \(h^1=0\), and evaluation map show
that it is pulled back from the base.  Hence \(\mathcal L^+\)
descends to a relative degree-\(b\) line bundle on the contracted
\(\mathbb P^1\)-bundle.  Fiberwise pullback identifies its sections
with those of \(\mathcal L^+\): a section is constant on every
contracted tail and is determined by its value on the marked component.
The corresponding \(H^0\) bundles agree by base change, so the
inclusion \(\mathcal L\to\mathcal L^+\) sends the original map
sections to sections of the descended bundle in families.  Restricted
to the marked component, this inclusion multiplies them by a section
vanishing to order \(e\) at the attachment of a tree of degree \(e\).
The original map sections have no common zero, so the new common zero
has exactly that multiplicity.  These constructions commute with base
change and produce the retained data in families.  Since the
contraction is an isomorphism near the marking, it preserves the
marked cotangent line.

\subsection{The weighted projective fiber}

Choose target coordinates with \(x=[1:0:\cdots:0]\).  Locally on
the base, identify the contracted bundle with a product
\(\mathbb P^1\)-bundle and the descended line bundle with
\(\mathcal O_{\mathbb P^1}(b)\).  Put the mark at infinity and
normalize the degree-\(b\) leading coefficient vector of the
polynomial sections to \((1,0,\ldots,0)\).  Translation of the
affine coordinate uniquely eliminates the coefficient of degree
\(b-1\) in the first polynomial.  After these normalizations, two
local choices differ only by scaling the affine coordinate.
Consequently the coefficient tuples glue as points of the quotient
stack for this residual scaling action.

The remaining coefficients have weights
\[
  \underbrace{1,\ldots,1}_{N},\quad
  \underbrace{2,\ldots,2}_{N+1},\quad\ldots,\quad
  \underbrace{b,\ldots,b}_{N+1}.
\]
They are not all zero.  Otherwise the polynomial tuple would be
\((z^b,0,\ldots,0)\), a constant map with all degree concentrated at
one basepoint.  A constant marked component has at least two distinct
attachments by stability, so cannot produce that tuple; a
nonconstant marked component cannot produce it either.

Let \(w_1,\ldots,w_M\) be the displayed list of weights and put
\[
 \mathcal W=[(\mathbb A^M\setminus\{0\})/\mathbb G_m],
 \qquad \lambda\cdot(x_1,\ldots,x_M)
       =(\lambda^{w_1}x_1,\ldots,\lambda^{w_M}x_M).
\]
This is the weighted projective stack in question. We have constructed
a proper morphism
\[
  \Phi:F\longrightarrow\mathcal W.
\]
The map is proper because \(F\) is proper and \(\mathcal W\) is
separated. For the reparametrization \(z\mapsto\lambda z\), the
normalized coefficient of \(z^{b-j}\) has weight \(-j\), and the
cotangent \(d(1/z)\) has weight \(-1\). Replacing \(\lambda\) by
\(\lambda^{-1}\) gives the displayed positive weights and identifies
the cotangent line with the pullback of
\(\mathcal O_{\mathcal W}(1)\), as in the one-marked weighted construction
of~\cite[proof of Lemma~3.3]{MustataMustata2007}.  Thus, with
\(\xi=c_1(\mathcal O_{\mathcal W}(1))\),
\[
  \psi=\Phi^*\xi.
\]

On the open substack \(U\), the tuple has no basepoints and recovers
the entire pointed map.  Its image is the locus of basepoint-free
tuples whose maps factor through \(\jmath(X)\) and have numerical
degree \(\beta\) there.  A residual scaling fixes the tuple exactly
when the associated reparametrization preserves the pointed map, so
the stabilizers are the automorphism groups of the pointed maps.  We
therefore have an identification of stacks onto this locus, with
generic stack degree one on every irreducible component of \(U\).

\subsection{The boundary does not contribute}

Consider a boundary stratum of \(F\).  Let \(d_0\) be the
anticanonical degree of the marked component and let \(u\ge1\) be
the number of trees directly attached to it.  Its image in
\(\mathcal W\) is determined by the marked map, the mark and
attachment points, and the degrees of those trees.  It does not
retain the maps on the trees: a tree of anticanonical degree \(d_i\)
is recorded only by its common-zero multiplicity \(ad_i\).

If \(d_0>0\), the marked map is free, since it passes through \(x\).
The space of its parametrized maps has dimension \(n+d_0\).
Adding \(u+1\) marked points, imposing evaluation \(x\), and dividing
by source reparametrization gives dimension
\[
  n+d_0+(u+1)-n-3=d_0+u-2.
\]
If \(d_0=0\), the map is the constant \(x\), stability gives \(u\ge2\),
and the remaining data have dimension
\(\dim\mathcal M_{0,u+1}=u-2\), the same formula.

Every attached tree contains a nonconstant rational component.
The pseudoindex assumption implies
\[
  d-d_0\ge2u,\qquad
  d_0+u-2\le d-u-2\le h-1.
\]
There are finitely many discrete boundary types.  We conclude that
\[
  \dim\Phi(F\setminus U)<h.
\]
This estimate concerns the image only; it makes no smoothness or
dimension assertion about the spaces of tail maps.

Write \(Z_\alpha\) for the closures in \(F\) of the irreducible
components of \(U\).  Chow localization and
\([F]^{\mathrm{vir}}|_U=[U]\) express
\[
  [F]^{\mathrm{vir}}-\sum_\alpha[Z_\alpha]
\]
as a class supported on \(F\setminus U\).  Its pushforward to
\(\mathcal W\) vanishes, since every \(h\)-dimensional boundary cycle
has image of smaller dimension.  The generic degree-one statement
on \(U\) therefore gives
\[
  \Phi_*[F]^{\mathrm{vir}}
    =\sum_\alpha[\overline{\Phi(U_\alpha)}].
\]
This is a nonempty effective sum of \(h\)-dimensional integral
stack cycles.  Virtual effectiveness has been established only
after pushforward; no effectiveness of \([F]^{\mathrm{vir}}\)
itself was assumed.

\subsection{Degrees in a weighted projective stack}

\begin{lemma}
\label{lem:weighted-degree-lower}
Let \(\mathcal P\) be a weighted projective stack with positive
integer weights at most \(b\), and let \(Z\subset\mathcal P\) be an
integral closed substack of dimension \(h\).  Then
\[
  \int_Z c_1(\mathcal O_{\mathcal P}(1))^h
  \ge b^{-(h+1)}.
\]
\end{lemma}

\begin{proof}
The affine cone of \(Z\) is defined by an ideal homogeneous for the
weighted grading.  Degenerate it to a monomial initial ideal.
This is a flat grading-preserving degeneration.  Removing the
origin and quotienting by the grading action of \(\mathbb G_m\)
gives a flat family of closed substacks of the fixed proper
weighted projective stack.  Its \(\mathcal O(1)\)-degree is constant.

The top-dimensional cycle of the monomial limit is a nonempty sum
of coordinate weighted projective substacks, with positive integer
multiplicities.  A component of dimension \(h\) has the form
\(\mathbb P(w_0,\ldots,w_h)\), with \(1\le w_j\le b\), and
\[
  \int_{\mathbb P(w_0,\ldots,w_h)}
       c_1(\mathcal O(1))^h
   =\frac1{w_0\cdots w_h}.
\]
For clarity, the power-coordinate map
\[
  \mathbb P^h\longrightarrow\mathbb P(w_0,\ldots,w_h),
  \qquad [x_0:\cdots:x_h]\longmapsto
          [x_0^{w_0}:\cdots:x_h^{w_h}]
\]
is the map of stacks specified by the line bundle
\(\mathcal O_{\mathbb P^h}(1)\) and these sections of its powers.
It is representable and finite, pulls back \(\mathcal O(1)\) to
\(\mathcal O_{\mathbb P^h}(1)\), and has stack degree
\(\prod_jw_j\).  On the dense tori the coarse degree is
\(\prod_jw_j/g\), where \(g=\gcd(w_0,\ldots,w_h)\); the target's
generic stabilizer has order \(g\), giving the stated stack degree.
This also covers \(h=0\), when the target is \(B\mu_{w_0}\).
The displayed integral follows by the projection formula.

Each component degree is at least \(b^{-(h+1)}\), so the same lower
bound holds for their nonempty positive sum and hence for \(Z\).
\end{proof}

\begin{proof}[Proof of Proposition~\ref{prop:descendant-bound}]
By the definition of the point descendant and by the projection formula,
\[
  \left\langle\tau_{d-2}(\mathrm{pt})\right\rangle_\beta
  =\int_{[F]^{\mathrm{vir}}}\psi^h
  =\int_{\Phi_*[F]^{\mathrm{vir}}}\xi^h.
\]
The last cycle is the nonempty effective sum obtained above.
All weights of \(\mathcal W\) are at most \(b\).
Lemma~\ref{lem:weighted-degree-lower} now gives the bound
\(b^{-(h+1)}=b^{-(d-1)}\).
\end{proof}

\subsection{Enough free degrees}

The next lemma lets us apply the descendant bound in a direction
avoiding any specified nonzero polynomial relation.  Write
\(N_1(X)_{\mathbb C}\) for the complex vector space of numerical
curve classes.

\begin{lemma}
\label{lem:free-dense}
The numerical degrees represented by nonconstant free rational maps
are Zariski dense in \(N_1(X)_{\mathbb C}\).  They are closed under
addition and under multiplication by positive integers.
\end{lemma}

\begin{proof}
We first prove closure under addition.  Take two free maps of degrees
\(\beta_1,\beta_2\).  Their evaluation images contain nonempty open
subsets of \(X\).  Deform them so that they meet at a common point
and glue there.  On the resulting two-component genus-zero curve
\(C\), let \(f:C\to X\) be the glued map and choose a smooth point
\(p\) of the first component.  Normalization gives
\[
  H^1(C,f^*T_X)=H^1(C,f^*T_X(-p))=0.
\]
Indeed the corresponding cohomology groups vanish on both
components, and sections on the second component evaluate
surjectively at the node.  Deformation theory therefore permits
a pointed smoothing of the node.  Semicontinuity of the second
vanishing shows that the smoothed map is free.  Its degree is
\(\beta_1+\beta_2\).

We next show that free degrees span the numerical vector space.
Smooth complex Fano manifolds are rationally chain connected by
\cite[Corollary~3.2]{Campana1992}, hence rationally connected by
\cite[Theorem~2.49]{Debarre2011}; in characteristic zero this is
equivalent to separable rational connectedness
\cite[Definition~2.20 and the following discussion]{Debarre2011}.
The rational-connectedness results are due to Campana and
Koll\'ar--Miyaoka--Mori~\cite{Campana1992,KollarMiyaokaMori1992}.
By \cite[Theorem~1.3]{TianZong2014}, the Chow group of one-cycles on \(X\)
is generated by rational curves.  In particular, rational curve
classes span \(N_1(X)_{\mathbb C}\).
A map is \emph{very free}
if its pulled-back tangent bundle is ample.  For a morphism
\(f_0:\mathbb P^1\to X\), comb smoothing attaches a nonempty finite
collection of very free rational maps as teeth and deforms the resulting
tree to a free rational map.  More precisely,
Tian--Zong~\cite[Proposition~2.4]{TianZong2014}, with handle
\(\mathbb P^1\) and line-bundle twist of degree \(-1\), gives a
smoothing \(f\) with \(H^1(\mathbb P^1,f^*T_X(-1))=0\).
Splitting on \(\mathbb P^1\) shows that \(f\) is free.  Its numerical
degree is the sum of the handle degree and the tooth degrees.
The tooth degrees are free, and their sum is free by closure under
addition.  Every rational degree is consequently a difference of
free degrees, which proves the spanning assertion.

Choose \(r=\dim N_1(X)_{\mathbb C}\) linearly independent free
degrees from the spanning set.  Every positive integral combination
of these degrees is free by the preceding addition argument.
Such combinations are Zariski dense: in the basis formed by the
chosen degrees, a polynomial vanishing on
\(\mathbb Z_{>0}^r\) vanishes by successive application of the
one-variable root bound.  Finally, covers of a free map remain
free, since pullback preserves the nonnegative summands of
its split tangent bundle.  This gives the assertion about positive
integer multiples directly as well.
\end{proof}

\section{An independent joint eigenvalue tuple}\label{sec:spectral}

We now combine the descendant lower bound with the recurrence to show
that the quantum divisors vary in all $r$ numerical directions.
Let
\[
 L=\overline{\mathbb C(q_1,\ldots,q_r)}.
\]
The commuting matrices $A_1,\ldots,A_r$ can be simultaneously put in
upper triangular form over $L$. A \emph{joint diagonal tuple} is the
$r$-tuple of diagonal entries at one position in such a form.

\begin{proposition}\label{prop:spectral}
At least one joint diagonal tuple $(\lambda_1,\ldots,\lambda_r)$ of the
quantum divisor matrices has coordinates algebraically independent over
$\mathbb C$.
\end{proposition}

We first normalize the descendant recurrence to obtain a uniform
exponential upper bound. We then show that a polynomial relation among
every joint tuple would force superexponential decay along a free degree,
contradicting Proposition~\ref{prop:descendant-bound}.

\subsection{Factorial normalization}

For $d_\beta\ge0$ and $0\le k\le n$, put
\[
 w_{\beta,k}=(d_\beta+k)!\,v_{\beta,k},
 \qquad w_\beta=\sum_{k=0}^n w_{\beta,k},
\]
where $v_{\beta,k}$ is the $H_k$ component. Put $w_\beta=0$ when
$d_\beta<0$. For a sequence $u$ indexed by $\mathbb Z^r$, let
$(E^\gamma u)_\beta=u_{\beta-\gamma}$ and set
$A_j(E)=\sum_\gamma A_{j,\gamma}E^\gamma$.
These are finite sums of shifts, and they commute.

If $A_{j,\gamma}$ sends an input component $H_l$ to $H_k$, then
$l=k-1+d_\gamma$. Thus
\[
 d_{\beta-\gamma}+l=d_\beta+k-1.
\]
Multiplying the $H_k$ component of \eqref{eq:recurrence} by
$(d_\beta+k-1)!$ gives, for $d_\beta>0$,
\begin{equation}\label{eq:normalized-recurrence}
 (A_j(E)w)_{\beta,k}
       =\frac{\beta_j}{d_\beta+k}\,w_{\beta,k}.
\end{equation}
Every nonzero input term has a nonnegative factorial argument. We do not
use this formula at degree zero.

Fix a norm on $H$. There is a constant $B>1$ such that
\begin{equation}\label{eq:exponential-bound}
 \|w_\beta\|\le B^{d_\beta+1}\qquad(d_\beta\ge0).
\end{equation}
To prove this, let $C_*=(\sum_jc_jD_j)\cup(-)$, and let $T_d$ multiply
$H_k$ by $1+k/d$ for positive integers $d$. Taking the linear
combination of \eqref{eq:normalized-recurrence} with coefficients $c_j$
yields
\[
 w_\beta=(I-T_{d_\beta}C_*)^{-1}T_{d_\beta}
              \sum_{\gamma>0}d_\gamma S_\gamma w_{\beta-\gamma}.
\]
Since $T_dC_*$ raises grading, the inverse $(I-T_dC_*)^{-1}$ is the
finite sum $\sum_{a=0}^n(T_dC_*)^a$. Its norm, and that of $T_d$, are uniformly
bounded for $d\ge1$. The sum uses finitely many positive shifts, each of
which strictly decreases $d$. Induction on the nonnegative integer $d$
proves \eqref{eq:exponential-bound}, taking $B$ large enough for the
finite set of coefficients and the initial value $w_0=n!\pt$.
All other degree-zero coefficients are zero.

\subsection{A relation would force too much decay}

\begin{proof}[Proof of Proposition~\ref{prop:spectral}]
Suppose, contrary to the proposition, that every joint diagonal tuple
is algebraically dependent over $\mathbb C$. There are finitely many
such tuples. Multiplying a nonzero polynomial relation for each of them
gives a nonzero $P\in\mathbb C[T_1,\ldots,T_r]$ vanishing on all of them.
If $m=\dim H$, simultaneous triangularization gives
$P(A_1,\ldots,A_r)^m=0$.
Substitute $(t^{c_j}q_j)_j$ and use \eqref{eq:homogeneity}. It follows
that $P(tA_1,\ldots,tA_r)^m=0$. The leading coefficient in $t$ gives
\begin{equation}\label{eq:homogeneous-relation}
 Q(A_1,\ldots,A_r)=0
\end{equation}
for a nonzero homogeneous polynomial $Q$: the $m$th power of the top
homogeneous part of $P$.

By Lemma~\ref{lem:free-dense}, choose a nonconstant free degree $\eta$
with $Q(\eta)\ne0$. Set $p_* =\eta/d_\eta$. Homogeneity gives
$Q(p_*)\ne0$. Choose a small bounded neighborhood of $p_*$ on whose
closure $|Q|$ is bounded below by a positive constant.
The substitution $q^\gamma\mapsto E^\gamma$ is a homomorphism from
Laurent polynomials to constant-coefficient shift operators. Applying it
entrywise to \eqref{eq:homogeneous-relation} gives
$Q(A_1(E),\ldots,A_r(E))=0$. Since these operators commute,
polynomial differencing gives
\begin{equation}\label{eq:polynomial-difference}
 Q(p)I=\sum_{j=1}^rB_j(p,E)(p_jI-A_j(E)).
\end{equation}
Here $p=(p_1,\ldots,p_r)$ consists of scalar parameters and every $B_j$
is a matrix polynomial in $p$ and finitely many shifts. This follows by
successively replacing the commuting variables in $Q(p)-Q(A(E))$.

Apply \eqref{eq:polynomial-difference} to $w$ at an index $\beta$,
with $p=\beta/d_\beta$ held fixed throughout this application.
For any shift $\delta$ that occurs and any $k$, the relevant component
of the rightmost factor is
\[
 \begin{aligned}
 ((p_jI-A_j(E))w)_{\beta-\delta,k}
 &=\left(\frac{\beta_j}{d_\beta}
    -\frac{\beta_j-\delta_j}{d_\beta-d_\delta+k}\right)
       w_{\beta-\delta,k}\\
 &=\frac{p_j(k-d_\delta)+\delta_j}{d_\beta-d_\delta+k}
       w_{\beta-\delta,k}.
 \end{aligned}
\]
When $p$ stays in the chosen neighborhood and $d_\beta$ is large,
the last scalar factor is uniformly $O(d_\beta^{-1})$: its numerator
is bounded, and its denominator differs from $d_\beta$ by a bounded amount.
The finitely many coefficient matrices in $B_j$ are uniformly bounded
there. Hence there are $C>0$ and a finite set of shifts $U$ such that
\begin{equation}\label{eq:small-shift-bound}
 \|w_\beta\|\le\frac C{d_\beta}
                     \max_{\delta\in U}\|w_{\beta-\delta}\|.
\end{equation}
If the right side of \eqref{eq:polynomial-difference} is the zero
operator, it directly implies $w_\beta=0$ in this neighborhood, which
already contradicts the descendant lower bound. We may therefore take
$U$ nonempty, including the zero shift if it occurs.

Take $\beta=s\eta$ for positive integers $s$. For sufficiently small
fixed $\epsilon>0$, every index reached by at most
$\lfloor\epsilon d_\beta\rfloor$ successive shifts from $U$ has degree
in $[d_\beta/2,2d_\beta]$ and direction in the chosen neighborhood.
Indeed each step changes both the vector and its degree by bounded
amounts; the total change is at most a constant times
$\epsilon d_\beta$. Choose $\epsilon$ first to keep the degree above
$d_\beta/2$, then to keep the normalized direction close to $p_*$.
For large $s$ all these indices also exceed the fixed lower threshold
for \eqref{eq:small-shift-bound}. Iteration, followed by
\eqref{eq:exponential-bound}, gives
\begin{equation}\label{eq:superexponential-bound}
 \|w_\beta\|\le
 \left(\frac{2C}{d_\beta}\right)^{\lfloor\epsilon d_\beta\rfloor}
 B^{2d_\beta+1}.
\end{equation}
This finite iteration does not require the shifts to decrease the
degree strictly, and it remains valid at indices with zero-extended
coefficients.

On the other hand, every positive multiple $s\eta$ is free: compose a
free map of degree $\eta$ with a degree-$s$ cover of $\mathbb P^1$.
By \eqref{eq:string-pairing} and
Proposition~\ref{prop:descendant-bound},
\[
 (w_\beta,1)
 \ge\frac{(d_\beta+n)!}{(a d_\beta)^{d_\beta-1}}
 \ge a\,d_\beta^{n+1}(ae)^{-d_\beta},
\]
where $a$ is fixed by the anticanonical embedding. The last inequality
uses $d!\ge(d/e)^d$. Pairing with $1$ is a bounded linear functional on
$H$, so this bounds $\|w_\beta\|$ below up to a fixed factor.
Its logarithm is bounded below by $-O(d_\beta)$, whereas the logarithm
of \eqref{eq:superexponential-bound} is
$-\epsilon d_\beta\log d_\beta+O(d_\beta)$. This contradiction proves
Proposition~\ref{prop:spectral}.
\end{proof}

\section{From divisor eigenvalues to independent curve degrees}
\label{sec:trace}

We now turn algebraic independence of a joint eigenvalue tuple into a
nonzero product of curve correspondences in independent degrees.  The
point requiring care is that a common triangular basis can depend on
the variables $q_1,\ldots,q_r$.  The following lemma accounts for the derivatives
of that basis and applies without a semisimplicity assumption.

For the Hochschild generalized-trace and inner-conjugation viewpoint,
see Khalkhali~\cite[Section~3.5, Examples~3.5.5--3.5.6 and Lemma~3.5.1]{Khalkhali2009}.
We give an elementary proof of the specific differential-form statement
needed here.

\begin{lemma}[An alternating trace form]
\label{lem:alternating-trace}
Let $L$ be a field extension of $\mathbb C$, let
$A_1,\ldots,A_r\in\operatorname{Mat}_m(L)$ commute pairwise, and let
$A_0\in\operatorname{Mat}_m(L)$ commute with each $A_j$, where $r\geq1$.
Use the K\"ahler differential $\dd$ of $L/\mathbb C$, and multiply
matrix-valued differential forms by matrix multiplication and exterior
multiplication.  Then
\begin{equation}
\label{eq:alternating-trace}
 \Theta(A_0;A_1,\ldots,A_r)
 :=\sum_{\pi\in\mathfrak S_r}\operatorname{sgn}(\pi)
 \operatorname{Tr}\bigl(A_0\,\dd A_{\pi(1)}\cdots \dd A_{\pi(r)}\bigr)
\end{equation}
is unchanged by replacing every $A_j$, including $A_0$, with
$GA_jG^{-1}$ for any $G\in\operatorname{GL}_m(L)$, and differentiating
the new entries.
\end{lemma}

\begin{proof}
Fix an arbitrary matrix-valued one-form $K$. For a scalar parameter $t$,
set
\[
 \begin{aligned}
 B_j(t)&=\dd A_j+t[K,A_j],\\
 F(t)&=\sum_{\pi\in\mathfrak S_r}\operatorname{sgn}(\pi)
       \operatorname{Tr}(A_0B_{\pi(1)}(t)\cdots B_{\pi(r)}(t)).
 \end{aligned}
\]
Thus $F(t)$ is a polynomial with values in differential forms.
Differentiating $[A_j,A_l]=0$ and using the Jacobi identity gives
\begin{equation}
\label{eq:commuting-differentials}
 [A_j,B_l(t)]=[A_l,B_j(t)].
\end{equation}
Here each commutator contains a degree-zero matrix, so it is the ordinary
commutator.

We prove that the derivative $F'(t)$ is
zero separately at each factor position.  Fix a position $s$ and a
permutation $\pi$, and abbreviate $C_h=B_{\pi(h)}(t)$ and $j=\pi(s)$.
Write
\[
 P=C_1\cdots C_{s-1},\qquad Q=C_{s+1}\cdots C_r.
\]
The contribution obtained by differentiating position $s$ is
\[
 T_{\pi,s}=\operatorname{Tr}(A_0PKA_jQ)
              -\operatorname{Tr}(A_0PA_jKQ).
\]
To compute it, move $A_j$ rightward through $Q$, cyclically from the
end of the trace to the beginning, through $A_0$, and then rightward
through $P$.  Expanding every commutator gives the exact identity
\begin{align}
 T_{\pi,s}
 &=\sum_{h<s}\operatorname{Tr}\bigl(
    A_0C_1\cdots C_{h-1}[A_j,C_h]
    C_{h+1}\cdots C_{s-1}KQ\bigr) \notag\\
 &\quad+\sum_{h>s}\operatorname{Tr}\bigl(
    A_0PKC_{s+1}\cdots C_{h-1}[A_j,C_h]
    C_{h+1}\cdots C_r\bigr).
 \label{eq:trace-telescoping}
\end{align}
Empty products have value $I$.  There is no term from crossing $A_0$
because $[A_j,A_0]=0$.  There is also no exterior sign in this identity:
the transported factor $A_j$ has degree zero, and no two one-forms
are exchanged.

For each fixed $h\ne s$, pair a permutation $\pi$ with the permutation
obtained by exchanging its labels at positions $s$ and $h$.
All factors in the corresponding summand of
\eqref{eq:trace-telescoping} remain the same except that
$[A_{\pi(s)},B_{\pi(h)}(t)]$ becomes
$[A_{\pi(h)},B_{\pi(s)}(t)]$.
These are equal by \eqref{eq:commuting-differentials}, while the
permutation sign changes.  Thus
\[
 \sum_{\pi\in\mathfrak S_r}\operatorname{sgn}(\pi)T_{\pi,s}=0
 \qquad\text{for every }s.
\]
The entire alternating trace polynomial is therefore constant in $t$.

Finally, set $K=G^{-1}\dd G$.  Direct differentiation yields
\[
 \dd(GA_jG^{-1})=G\bigl(\dd A_j+[K,A_j]\bigr)G^{-1}.
\]
In a product of these forms the degree-zero matrices $G^{-1}G$ cancel,
and the trace is invariant under conjugation.  Constancy between $t=0$
and $t=1$ proves the assertion.  Notice that the argument never
differentiates $A_0$ and does not require it to commute with $\dd A_j$.
\end{proof}

\begin{lemma}[A nonzero trace from one eigenvalue tuple]
\label{lem:spectral-trace}
Let $L$ be an algebraically closed extension of $\mathbb C$ and let
$A_1,\ldots,A_r\in\operatorname{Mat}_m(L)$ commute pairwise.
Suppose a joint diagonal tuple $(\lambda_1,\ldots,\lambda_r)$ in a
common upper triangular basis is algebraically independent over
$\mathbb C$.  There is a matrix $A_0$, polynomial over $L$ in the
$A_j$, for which \eqref{eq:alternating-trace} is nonzero.
\end{lemma}

\begin{proof}
Let $\Lambda$ be the finite set of distinct joint diagonal tuples,
and write $\lambda=(\lambda_1,\ldots,\lambda_r)$ for the chosen tuple.
For each $\mu\in\Lambda\setminus\{\lambda\}$ choose an index $j(\mu)$
with $\lambda_{j(\mu)}\ne\mu_{j(\mu)}$.  The polynomial
\[
 p(T_1,\ldots,T_r)
 =\prod_{\mu\in\Lambda\setminus\{\lambda\}}
   \frac{T_{j(\mu)}-\mu_{j(\mu)}}
        {\lambda_{j(\mu)}-\mu_{j(\mu)}}
\]
has value $1$ at $\lambda$ and $0$ at every other tuple.  Set
$A_0=p(A_1,\ldots,A_r)$.  This matrix commutes with every $A_j$.
We only use its prescribed diagonal entries; it need not be an
idempotent on a generalized eigenspace.

By Lemma~\ref{lem:alternating-trace}, the form can be computed in the
common upper triangular basis.  Derivatives and products remain upper
triangular, and their diagonal entries multiply.  If the chosen tuple
occurs with multiplicity $m_\lambda$, it follows that
\[
 \Theta(A_0;A_1,\ldots,A_r)
 =r!m_\lambda\,\dd\lambda_1\wedge\cdots\wedge \dd\lambda_r.
\]
The alternating sign and the sign from reordering the scalar one-forms
cancel for each permutation.  Algebraic independence in characteristic
zero makes the displayed wedge nonzero; the integer $r!m_\lambda$ is
also nonzero.  This proves the lemma.
\end{proof}

We apply these lemmas over
$L=\overline{\mathbb C(q_1,\ldots,q_r)}$ to the quantum divisor
operators.  Common triangularization is available for any finite
commuting collection over an algebraically closed field.  The
differential extends to $L$, since the algebraic extension is
separable.

\begin{proposition}[Independent curve correspondences]
\label{prop:independent-correspondences}
Suppose that the quantum divisor operators $A_1(q),\ldots,A_r(q)$
have an algebraically independent joint eigenvalue tuple over
$\mathbb C$.  There are linearly independent numerical degree vectors
$\gamma_1,\ldots,\gamma_r$ of positive genus-zero stable maps such that
\begin{equation}
\label{eq:nonzero-independent-product}
 S_{\gamma_1}\cdots S_{\gamma_r}\ne0,
 \qquad
 \sum_{j=1}^r(d_{\gamma_j}-1)\le n.
\end{equation}
Consequently $r(\iota-1)\le n$.
\end{proposition}

\begin{proof}
Choose $A_0$ by Lemma~\ref{lem:spectral-trace}. In the original fixed
cohomology basis, the classical terms of $A_j$ are constant, and hence
\[
 \dd A_j=\sum_{\gamma>0}q^\gamma\gamma_jS_\gamma\,\omega_\gamma,
 \qquad
 \omega_\gamma=\sum_{l=1}^r\gamma_l\frac{\dd q_l}{q_l}.
\]
These sums are finite. Each ordered summand in the expansion of
\eqref{eq:alternating-trace} is a scalar multiple of
\[
 \operatorname{Tr}(A_0S_{\gamma_1}\cdots S_{\gamma_r})\,
 \omega_{\gamma_1}\wedge\cdots\wedge\omega_{\gamma_r}.
\]
Since $\Theta\ne0$, at least one of these ordered summands is nonzero.
Its trace factor and its wedge factor are therefore both nonzero.
The forms $\dd q_l/q_l$ are an $L$-basis of the differentials of
$L/\mathbb C$, so the nonzero wedge gives linear independence of
the degree vectors. The nonzero trace gives the operator product in
\eqref{eq:nonzero-independent-product}, in the order selected by the
summand.

The operator $S_\gamma$ maps $H_k$ into $H_{k+1-d_\gamma}$. Thus this
product lowers the codimension grading by $\sum_j(d_{\gamma_j}-1)$.
The available gradings run only from $0$ through $n$, proving the upper
bound in \eqref{eq:nonzero-independent-product}. Every positive
stable-map degree has anticanonical degree at least $\iota$, so
\[
 r(\iota-1)\le\sum_{j=1}^r(d_{\gamma_j}-1)\le n.
\]
\end{proof}

Proposition~\ref{prop:spectral} supplies the algebraically independent
joint tuple required in Proposition~\ref{prop:independent-correspondences}.
This proves the inequality in Theorem~\ref{thm:main}; the equality case
is considered next.

\section{The equality case}
\label{sec:equality}

We show that equality in the grading bound forces the projective-space
product.  The nonzero product from
Proposition~\ref{prop:independent-correspondences} first gives a chain
of actual minimum-degree rational curves.  A dimension estimate then
turns the corresponding families into covering families.

We use $\operatorname{RatCurves}^{\mathrm n}(X)$ for the normalized
parameter space of rational curves, obtained from maps
$\mathbb P^1\to X$ birational onto their images by reparametrization.
In this section a family $V$ is an irreducible component of that space.
It is \emph{unsplit} if it is proper, and \emph{covering} if its
universal evaluation map dominates $X$.  Its anticanonical degree is
$-K_X\cdot V$, the common degree of its members.  For unsplit families
$V^1,\ldots,V^k$ and a point $x\in X$, write
$\operatorname{Locus}(V^1,\ldots,V^k)_x$ for the set of endpoints $y$
of chains $C_1,\ldots,C_k$ with $C_j\in V^j$, $x\in C_1$,
$C_j\cap C_{j+1}\ne\varnothing$, and $y\in C_k$.

We will use the following two established results, with their family
hypotheses stated explicitly.

\begin{theorem}[The ordered-chain bound]
\label{thm:ordered-chain-bound}
Let $Y$ be a smooth complex projective variety.  Let
$V^1,\ldots,V^k$ be proper irreducible components of
$\operatorname{RatCurves}^{\mathrm n}(Y)$ whose numerical classes are
linearly independent.  For every $y\in Y$, the ordered endpoint locus
is either empty or satisfies
\[
 \dim\operatorname{Locus}(V^1,\ldots,V^k)_y
 \ge\sum_{j=1}^k(-K_Y\cdot V^j-1).
\]
\end{theorem}

This is the theorem of Bonavero--Casagrande--Debarre--Druel
\cite[Theorem~5.2]{BCDD2003}, with the irreducible-component hypothesis
specified in their erratum~\cite{BCDDErratum}.  Taking arbitrary
proper subfamilies would not suffice.

\begin{theorem}[Occhetta's product characterization]
\label{thm:occhetta-product}
Let $Y$ be a smooth complex projective variety of dimension $n$.
Suppose $Y$ has numerically independent unsplit covering families
$V^1,\ldots,V^k$ with anticanonical degrees $n_1+1,\ldots,n_k+1$,
where each $n_j$ is positive and $\sum_jn_j=n$.  Then
$Y\cong\mathbb P^{n_1}\times\cdots\times\mathbb P^{n_k}$.
\end{theorem}

We use Theorem~\ref{thm:occhetta-product} in the form of
\cite[Theorem~1.1]{Occhetta2006}.  The next lemma verifies that the curve
families needed for these results are supplied by our correspondences.

\begin{lemma}[Minimum-degree incidence]
\label{lem:minimum-degree-incidence}
Let $\gamma_1,\ldots,\gamma_k$ be linearly independent numerical
stable-map degrees on $X$, each of anticanonical degree $\iota$.
If $S_{\gamma_1}\cdots S_{\gamma_k}\ne0$, there exist numerically
independent unsplit families $V^1,\ldots,V^k$, each of degree $\iota$,
and a nonempty chain in some ordering of those families.
\end{lemma}

\begin{proof}
For a numerical degree $\gamma$, let
$\overline{\mathcal M}_{0,2}(X,\gamma)$ denote the finite disjoint
union of stable-map stacks over actual curve classes of that degree.
Let $Z_\gamma\subset X\times X$ be the union of the images of their
two-evaluation maps.
It is closed by properness.  The correspondence inducing $S_\gamma$
is the pushforward of its virtual class and is supported on
$Z_\gamma$.  The composition of these correspondences is computed
on $X^{k+1}$ by imposing the consecutive-pair incidence conditions
and projecting to the two endpoints.  Refined intersection with the
diagonals is supported on the incidence fiber product.  If that fiber
product were empty, the resulting correspondence, and hence its
action on cohomology, would be zero.  Thus the nonzero product gives
a sequence of two-pointed stable maps with matching evaluations.
Choose their order according to composition, reversing the written
list of degrees if necessary.  This support argument does not require
effectivity of the virtual classes.

Every nonconstant component of a genus-zero stable map has an
irreducible rational image of anticanonical degree at least $\iota$.
Its contribution to the total degree is that image degree multiplied
by its covering multiplicity.  A stable map of total degree $\iota$
therefore has exactly one nonconstant component, and that component
is birational onto its image.  The images in our sequence are thus
irreducible rational curves $C_1,\ldots,C_k$ in the specified
numerical classes, with consecutive intersections nonempty.
Contracted components carrying markings do not affect this conclusion.

Choose an irreducible component $V^j$ of
$\operatorname{RatCurves}^{\mathrm n}(X)$ containing $C_j$.
Its members retain degree $\iota$.  To check properness, take stable
limits after forgetting markings.  Such a limit still has just one
nonconstant component, birational onto its image, by the same
degree argument.  It has no contracted components: any contracted
subtree would have a leaf of valency one and no markings, contrary
to stability.  Hence every unmarked limit has smooth domain
$\mathbb P^1$ and is birational onto its image.  Smooth-domain
birational stable maps represent the unparameterized birational-map
quotient underlying $\operatorname{RatCurves}(X)$.  The full
degree-$\iota$ stable-map locus is proper and has no boundary, so the
corresponding full irreducible component of this quotient is proper;
its finite normalization $V^j$ is proper as well.
These are actual proper irreducible components, as required by
Theorem~\ref{thm:ordered-chain-bound}.  Their numerical classes are
the independent classes $\gamma_j$, and the curves already selected
give the required chain.
\end{proof}

\begin{proposition}[Rigidity at equality]
\label{prop:equality-rigidity}
Suppose $\iota\ge2$, $r(\iota-1)=n$, and the independent nonzero
product in \eqref{eq:nonzero-independent-product} exists.  Then
\[
 X\cong(\mathbb P^{\iota-1})^r.
\]
\end{proposition}

\begin{proof}
Since each $d_{\gamma_j}\ge\iota$, equality in the grading bound
forces $d_{\gamma_j}=\iota$ for every $j$.
Lemma~\ref{lem:minimum-degree-incidence} supplies numerically
independent proper components $V^1,\ldots,V^r$ and curves
$C_1,\ldots,C_r$ forming a chain in that order.  Choose $x\in C_1$.
Theorem~\ref{thm:ordered-chain-bound} gives
\[
 \dim\operatorname{Locus}(V^1,\ldots,V^r)_x
 \ge r(\iota-1)=n.
\]
This locus is contained in the image of the universal evaluation of
$V^r$, so $V^r$ is covering.  Since $V^r$ is proper, its evaluation
image is closed; consequently every point of $X$ lies on a member
of $V^r$.

If $r>1$, choose a member of $V^r$ through a point of $C_1$ and
follow it by $C_1,\ldots,C_{r-1}$.  This is a chain in the order
$V^r,V^1,\ldots,V^{r-1}$.  Applying the same dimension bound proves
that $V^{r-1}$ is covering.  Repeating this cyclic rotation makes
each family the last member of a nonempty ordered chain, so every
$V^j$ is covering.  For $r=1$ the first application already suffices.

All hypotheses of Theorem~\ref{thm:occhetta-product} now hold with
$n_j=\iota-1$: the families are unsplit and covering, their
numerical classes are independent, their degrees are $\iota$, and
their contributions $n_j$ sum to $n$.  The resulting product is
$(\mathbb P^{\iota-1})^r$.
\end{proof}

Conversely, $(\mathbb P^{\iota-1})^r$ has dimension $r(\iota-1)$
and Picard number $r$.  Its anticanonical divisor has degree $\iota$
on a line in one factor, and degree at least $\iota$ on every
irreducible rational curve.  Its pseudoindex is therefore $\iota$,
and equality holds.

\bibliographystyle{plain}
\bibliography{references}
\end{document}
